\section{Source Section II.B: Gradient Entropy and Energy}
\label{sec:gradient-entropy-energy}

\sourceeqrange{Eqs. (2.12)--(2.20)}

\subsection{Purpose and Scope}

This is the companion's first worked derivation section.  It follows the
source equation order in Section II.B, while reconstructing the typed calculation
behind the formulas in the companion's own notation.  The source paper should be
read alongside this guide.

The goal here is narrow:

\begin{itemize}
\item define the entropy and internal-energy functionals with square-gradient
  terms;
\item keep entropy density, entropy per particle, internal-energy density, and
  gradient contributions distinct;
\item derive why the coefficient appearing in the equilibrium gradient free
  energy is
  \[
    M=CT+K;
  \]
\item reconstruct the generalized chemical-potential variation at fixed
  gradient-augmented internal energy;
\item isolate the boundary contribution generated by integration by parts.
\end{itemize}

The reversible stress tensor, hydrodynamic equations, Appendix-B scaled
equations, and numerical assumptions are deferred to later sections.  In
particular, the Appendix-B B3 branch issue is not a Section II.B result.

\subsection{Typed Objects and Regularity Assumptions}

Let \(\Omega\subset\mathbb{R}^d\) be the fluid container and let
\(\partial\Omega\) be a piecewise smooth boundary with outward unit normal
\(\nu\).  The source uses \(d=1,2,3\) in different contexts; the formulas in
this section are local in \(\mathbb{R}^d\).

In this section, subscripts such as \(C_n\), \(K_n\), and \(M_n|_T\) are
thermodynamic partial derivatives with respect to the scalar argument \(n\).
Spatial derivatives are written with \(\grad\), \(\grad\cdot\), or component
derivatives.  The typed fields are as follows.

\begin{center}
\begin{tabular}{@{}lll@{}}
\toprule
Symbol & Type & Role in this section \\
\midrule
\(n:\Omega\to\mathbb{R}_{>0}\) & scalar field & number-density field \\
\(e:\Omega\to\mathbb{R}\) & scalar density & internal-energy density without the gradient term \\
\(\hat e:\Omega\to\mathbb{R}\) & scalar density & internal-energy density including the gradient term \\
\(s=s(n,e)\) & scalar per particle & entropy per particle \\
\(\hat S\) & scalar density & entropy density including the gradient term \\
\(T=T(n,e)\) & positive scalar field & local temperature defined by the entropy derivative \\
\(\mu=\mu(n,e)\) & scalar per particle & ordinary chemical potential \\
\(\hat\mu\) & scalar per particle & generalized chemical potential including gradients \\
\(C=C(n)\) & coefficient field & entropy-gradient coefficient \\
\(K=K(n)\) & coefficient field & energy-gradient coefficient \\
\(M=M(n,T)\) & coefficient field & equilibrium gradient free-energy coefficient \\
\(\delta n,\delta\hat e\) & variations & independent variations used after choosing \((n,\hat e)\) \\
\bottomrule
\end{tabular}
\end{center}

For a classical derivation, assume \(n\in C^2(\Omega)\),
\(\hat e,T,C,K,M\in C^1(\Omega)\), \(T>0\), and variations
\(\delta n,\delta\hat e\in C^1(\overline{\Omega})\).  A weak version would
replace these by Sobolev hypotheses with a trace map; this companion records
the classical calculation first.


\paragraph{Definition and relation status.}
The following status ledger fixes how the companion uses the Section II.B
objects.  It prevents a definition, a local thermodynamic identity, and a later
boundary or hydrodynamic statement from being read as the same kind of claim.
\begin{description}[leftmargin=2.8cm,style=nextline]
\item[\(S_b,E_b,\hat S,\hat e\)] source-stated model functionals and
pointwise densities for the square-gradient entropy/energy branch.
\item[\(\dd e=T\dd(ns)+\mu\dd n\)] Local equilibrium thermodynamic relation;
not an integrated balance law.
\item[\(1/T=n(\partial s/\partial e)_n\)] derived coefficient identity from
the entropy-density relation; this fixes the factor \(n\) in Eq.~(2.17).
\item[\(M=CT+K\)] Derived local coefficient identity, not a boundary condition
or stress law.
\item[\(\hat\mu\)] Variational derivative at fixed \(\hat e\), not yet a force
density.
\item[Eq.~\eqref{eq:iib-entropy-variation-boundary}] Boundary term from
integration by parts; a fixed-trace, natural-wall, or wall-functional branch is
needed before any global statement.
\end{description}

\subsection{Functional Definitions}

The source functionals in Eqs.~(2.12)--(2.15) are most transparent when read as
two levels of density.

The entropy functional is
\begin{equation}
  S_b[n,e]
  =
  \int_{\Omega}
  \left(
    n\,s(n,e)-\frac{1}{2}C(n)|\grad n|^2
  \right)\dd x .\sourceeqbadge{2.12}
  \label{eq:iib-entropy-functional}
\end{equation}
Here \(n\,s(n,e)\) is an entropy density, while
\(-C|\grad n|^2/2\) is a square-gradient entropy-density correction.

The internal-energy functional is
\begin{equation}
  E_b[n,e]
  =
  \int_{\Omega}
  \left(
    e+\frac{1}{2}K(n)|\grad n|^2
  \right)\dd x .\sourceeqbadge{2.13}
  \label{eq:iib-energy-functional}
\end{equation}
The associated pointwise gradient-augmented densities are
\begin{equation}
  \hat S
  =
  n\,s(n,e)-\frac{1}{2}C(n)|\grad n|^2,
  \qquad
  \hat e
  =
  e+\frac{1}{2}K(n)|\grad n|^2 .\sourceeqbadge{2.14--2.15}
  \label{eq:iib-hatted-densities}
\end{equation}
Thus \(e\) and \(\hat e\) are different density variables.  When \(S_b\) is
viewed as a functional of \((n,\hat e)\), the bare internal-energy density is
\begin{equation}
  e
  =
  \hat e-\frac{1}{2}K(n)|\grad n|^2 .
  \label{eq:iib-e-from-ehat}
\end{equation}

\subsection{Local Temperature and Basic Thermodynamic Differential}

The local temperature in \eqsrc{2.17} is a coefficient statement for the
nongradient entropy density
\begin{equation}
  \sigma(n,e)=n\,s(n,e).
  \label{eq:iib-entropy-density-sigma}
\end{equation}
The local equilibrium premise is the pointwise fundamental relation for the
internal-energy density,
\begin{equation}
  \dd e=T\,\dd\sigma+\mu\,\dd n .
  \label{eq:iib-local-fundamental-relation}
\end{equation}
Solving \eqref{eq:iib-local-fundamental-relation} for \(\dd\sigma\) gives
\begin{equation}
  \dd\sigma=\frac{1}{T}\dd e-\frac{\mu}{T}\dd n .
  \label{eq:iib-local-fundamental-solved}
\end{equation}
Independently, the chain rule for \(\sigma=n s(n,e)\) gives
\begin{align}
  \dd\sigma
  &=\dd(ns) \\
  &=s\,\dd n+n\,\dd s \\
  &=\left[s+n\left(\frac{\partial s}{\partial n}\right)_e\right]\dd n
    +n\left(\frac{\partial s}{\partial e}\right)_n\dd e .
  \label{eq:iib-sigma-chain-rule}
\end{align}
Comparing the coefficients of \(\dd e\) in
\eqref{eq:iib-local-fundamental-solved} and
\eqref{eq:iib-sigma-chain-rule} gives
\begin{equation}
  \frac{1}{T}
  =
  n\left(\frac{\partial s}{\partial e}\right)_n .\sourceeqbadge{2.17}
  \label{eq:iib-temperature-definition}
\end{equation}
This is why the factor \(n\) appears in the source equation: \(s\) is entropy
per particle, while \(\sigma=ns\) is the entropy density.

The same coefficient can be checked directly from the Section II.A van der
Waals formulas.  At fixed \(n\),
\begin{align}
  \left(\frac{\partial e}{\partial T}\right)_n
  &=\frac{d}{2}n k_B,\\
  \left(\frac{\partial s}{\partial T}\right)_n
  &=\frac{d k_B}{2T}.
  \label{eq:iib-section-iia-fixed-n-derivatives}
\end{align}
Therefore
\begin{align}
  \left(\frac{\partial s}{\partial e}\right)_n
  &=
  \frac{(\partial s/\partial T)_n}{(\partial e/\partial T)_n} \\
  &=
  \frac{d k_B/(2T)}{d n k_B/2}
  =
  \frac{1}{nT}.
  \label{eq:iib-section-iia-temperature-check}
\end{align}
Multiplying by \(n\) again gives \eqref{eq:iib-temperature-definition}.  This
source-specific route is only a consistency check; the general Section II.B
derivation uses the local equilibrium relation
\eqref{eq:iib-local-fundamental-relation}.

For a perturbation \(e_\varepsilon=e+\varepsilon\xi\) at fixed \(n\), the
first variation form is now a consequence of
\eqref{eq:iib-temperature-definition}:
\begin{align}
  \left.
  \frac{\dd}{\dd\varepsilon}
  \sigma(n,e+\varepsilon\xi)
  \right|_{\varepsilon=0}
  &=
  \left(\frac{\partial\sigma}{\partial e}\right)_n\xi \\
  &=
  n\left(\frac{\partial s}{\partial e}\right)_n\xi
  =
  \frac{1}{T}\xi .
  \label{eq:iib-temperature-first-variation-check}
\end{align}

The ordinary chemical potential \(\mu\) is used through the local differential
identity
\begin{equation}
  \delta(n s)
  =
  \frac{1}{T}\,\delta e
  -
  \frac{\mu}{T}\,\delta n .
  \label{eq:iib-local-thermo-differential}
\end{equation}
Equation \eqref{eq:iib-local-thermo-differential} is a pointwise local
equilibrium identity for the nongradient entropy density.  It is not an
integrated balance law and not a boundary condition.

\subsection{Variation at Fixed \texorpdfstring{\(\hat e\)}{ehat}}

The generalized chemical potential in \eqsrc{2.18} is defined by varying
\(S_b\) with respect to \(n\) while holding \(\hat e\), not \(e\), fixed.
This is the main bookkeeping point in Section II.B.  The chain of dependence is
\[
  (n,\hat e) \longmapsto e=\hat e-\frac12K(n)|\grad n|^2
  \longmapsto n s(n,e)
  \longmapsto \hat S=n s(n,e)-\frac12 C(n)|\grad n|^2 .
\]
Thus a variation of \(n\) at fixed \(\hat e\) still changes the bare energy
\(e\), because \(e\) contains the negative of the \(K\)-gradient correction.
Use the perturbation
\[
  n_\varepsilon=n+\varepsilon\eta,
  \qquad
  \hat e_\varepsilon=\hat e+\varepsilon\xi,
\]
where \(\eta=\delta n\) and \(\xi=\delta\hat e\).  To first order,
\begin{align}
  K(n_\varepsilon)|\grad n_\varepsilon|^2
  &=K|\grad n|^2
    +\varepsilon\left(K_n\eta|\grad n|^2
    +2K\grad n\cdot\grad\eta\right)+O(\varepsilon^2),
  \label{eq:iib-K-first-order-expansion}\\
  C(n_\varepsilon)|\grad n_\varepsilon|^2
  &=C|\grad n|^2
    +\varepsilon\left(C_n\eta|\grad n|^2
    +2C\grad n\cdot\grad\eta\right)+O(\varepsilon^2).
  \label{eq:iib-C-first-order-expansion}
\end{align}
Consequently the variation of \eqref{eq:iib-e-from-ehat} is
\begin{equation}
  \delta e
  =
  \delta\hat e
  -
  \frac{1}{2}K_n|\grad n|^2\delta n
  -
  K\,\grad n\cdot\grad\delta n ,
  \label{eq:iib-delta-e}
\end{equation}
where \(K_n=\mathrm{d}K/\mathrm{d}n\).  The signs in
\eqref{eq:iib-delta-e} come from differentiating
\(-K(n)|\grad n|^2/2\), not from an entropy convention.  In component form,
\[
  \delta\left[-\frac12K(n)\partial_i n\partial_i n\right]
  =-\frac12K_n\delta n\,\partial_i n\partial_i n
   -K\partial_i n\,\partial_i\delta n,
\]
with the repeated index summed.  The local thermodynamic differential
then contributes
\begin{equation}
  \delta(ns)
  =
  \frac{1}{T}\xi
  -\frac{\mu}{T}\eta
  -\frac{K_n}{2T}|\grad n|^2\eta
  -\frac{K}{T}\grad n\cdot\grad\eta .
  \label{eq:iib-local-entropy-first-variation-expanded}
\end{equation}
The variation of the entropy-gradient term is
\begin{equation}
  \delta\!\left[-\frac{1}{2}C|\grad n|^2\right]
  =
  -\frac{1}{2}C_n|\grad n|^2\delta n
  -
  C\,\grad n\cdot\grad\delta n .
  \label{eq:iib-delta-gradient-entropy}
\end{equation}

Insert \eqref{eq:iib-delta-e} and
\eqref{eq:iib-delta-gradient-entropy} into the variation of
\eqref{eq:iib-entropy-functional}, with \(\delta n=\eta\) and
\(\delta\hat e=\xi\).  A useful way to check the calculation by hand is to
sort terms by the object multiplying the variation:
\[
  \delta S_b
  =\int_\Omega
   \left(
     [\delta\hat e] + [\delta n] + [\grad\delta n]
   \right)\dd x,
\]
where the bracket labels denote type classes, not new source notation.  The
four coefficient groups before integration by parts
are
\begin{align}
  [\delta\hat e] &: \quad \frac{1}{T}, \\
  [\delta n] &: \quad -\frac{\mu}{T}
    -\frac{1}{2}\left(\frac{K_n}{T}+C_n\right)|\grad n|^2, \\
  [\grad\delta n] &: \quad -\left(\frac{K}{T}+C\right)\grad n, \\
  [\partial\Omega] &: \quad \text{not yet present; it appears only after integration by parts.}
  \label{eq:iib-pre-ibp-coefficient-ledger}
\end{align}
This ledger is a companion audit device: it records the type of each term before
any source-compatible abbreviation is introduced.  Before any integration by
parts, the first variation is
\begin{align}
  \delta S_b
  =\int_\Omega
  \Bigg[
    \frac{1}{T}\delta\hat e
    -\frac{\mu}{T}\delta n
    -\frac{1}{2}\left(\frac{K_n}{T}+C_n\right)|\grad n|^2\delta n
    -\left(\frac{K}{T}+C\right)\grad n\cdot\grad\delta n
  \Bigg]\dd x .
  \label{eq:iib-first-variation-before-M}
\end{align}
The terms multiplying \(\grad n\cdot\grad\delta n\) combine as
\begin{equation}
  -\left(\frac{K}{T}+C\right)\grad n\cdot\grad\delta n .
  \label{eq:iib-gradient-variation-combination}
\end{equation}
This is the calculation point where the coefficient relation enters.

\subsection{Why \texorpdfstring{\(M=CT+K\)}{M=CT+K}}

The equilibrium square-gradient free-energy coefficient is the coefficient of
\(|\grad n|^2/2\) in the Helmholtz free-energy density \(e-Tns\) augmented by
the two gradient corrections.  Combining the energy-gradient contribution
\(K|\grad n|^2/2\) with the entropy-gradient contribution after multiplication
by \(T\) gives
\begin{equation}
  \frac{1}{2}K|\grad n|^2
  +
  \frac{1}{2}TC|\grad n|^2
  =
  \frac{1}{2}(CT+K)|\grad n|^2 .
  \label{eq:iib-M-relation-derivation}
\end{equation}
Therefore
\begin{equation}
  M=CT+K .\sourceeqbadge{2.16}
  \label{eq:iib-M-relation}
\end{equation}
Equivalently,
\[
  \frac{M}{T}=C+\frac{K}{T}.
\]
The coefficient of \(\grad n\cdot\grad\delta n\) in
\eqref{eq:iib-gradient-variation-combination} is therefore
\[
  -\frac{M}{T}\,\grad n\cdot\grad\delta n .
\]

This coefficient identity is local and algebraic.  It does not by itself prove
any hydrodynamic equation, boundary condition, or numerical branch.

\subsection{Bulk Part of the Generalized Chemical Potential}

The coefficient of \(|\grad n|^2\delta n/2\) also becomes the
fixed-temperature derivative used by the source.  The distinction is operational:
\(M_n|_T\) differentiates the thermodynamic function \(M(n,T)\) while holding
its second argument fixed, whereas a spatial derivative of the field
\(M(n(x),T(x))\) would include both \(\grad n\) and \(\grad T\).  In formulas,
\begin{equation}
  \grad M=M_n|_T\,\grad n+M_T|_n\,\grad T,
  \label{eq:iib-spatial-versus-thermo-derivative}
\end{equation}
but only the first coefficient enters the fixed-temperature term below:
\begin{equation}
  \frac{K_n}{T}+C_n=\frac{1}{T}(K_n+TC_n)=\frac{M_n|_T}{T}.
  \label{eq:iib-Mn-fixed-temperature-step}
\end{equation}
This step holds \(T\) fixed while differentiating \(M(n,T)=C(n)T+K(n)\) with
respect to \(n\).  It should not be confused with the spatial derivative
\(\grad M\), where both \(n(x)\) and \(T(x)\) may vary in space.  The
fixed-variable distinction is
\begin{align}
  M_n|_T\,\delta n
  &= (TC_n+K_n)\,\delta n, \
  \grad M
  &= (TC_n+K_n)\grad n+C\grad T.
  \label{eq:iib-Mn-versus-gradM-linebyline}
\end{align}
The first line is the thermodynamic variation at fixed \(T\) and belongs in
the fixed-\(\hat e\) coefficient of \(\delta n\).  The second line is the
spatial derivative when \(M=CT+K\) and appears later in the hydrodynamic
coefficient identity.

Using \(M=CT+K\), the pre-integration-by-parts variation can be written in the
source-compatible form
\begin{equation}
  \delta S_b
  =
  \int_{\Omega}
  \left[
    \frac{1}{T}\delta\hat e
    -
    \frac{\mu}{T}\delta n
    -
    \frac{M_n|_T}{2T}|\grad n|^2\delta n
    -
    \frac{M}{T}\grad n\cdot\grad\delta n
  \right]\dd x .
  \label{eq:iib-pre-ibp-variation}
\end{equation}
The derivative \(M_n|_T=(\partial M/\partial n)_T\) is the derivative used in
\eqsrc{2.18}.  The notation matters because \(M=M(n,T)\), while \(T\) is
also a field in nonequilibrium settings.

Let
\[
  A=\frac{M}{T}.
\]
The integration-by-parts identity for the last term starts with the pointwise
product-rule calculation
\begin{align}
  -A\grad n\cdot\grad\delta n
  &=-\grad\cdot(A\grad n\,\delta n)
    +\grad\cdot(A\grad n)\,\delta n,
  \label{eq:iib-pointwise-product-rule}\\
  -\int_{\Omega}A\,\grad n\cdot\grad\delta n\,\dd x
  &=
  \int_{\Omega}\grad\cdot(A\grad n)\,\delta n\,\dd x
  -
  \int_{\partial\Omega}A\,(\nu\cdot\grad n)\,\delta n\,\dd a .
  \label{eq:iib-ibp-identity}
\end{align}
After this step, the boundary contribution is already fixed by the outward-normal
convention:
\begin{equation}
  -\int_{\partial\Omega}\frac{M}{T}(\nu\cdot\grad n)\delta n\,\dd a .
  \label{eq:iib-boundary-term-linebyline}
\end{equation}
This sign is not adjustable without changing the convention for \(\nu\).  After
separating this boundary term, the bulk coefficient of \(\delta n\) is
\begin{align}
  -\frac{\mu}{T}
  -\frac{M_n|_T}{2T}|\grad n|^2
  +\grad\cdot\left(\frac{M}{T}\grad n\right)
  &=-\frac{1}{T}
  \left[
    \mu
    -
    T\grad\cdot\left(\frac{M}{T}\grad n\right)
    +
    \frac{M_n|_T}{2}|\grad n|^2
  \right].
  \label{eq:iib-bulk-coefficient-expanded}
\end{align}
This identifies the generalized chemical potential
\begin{equation}
  \hat\mu
  =
  \mu
  -
  T\grad\cdot\left(\frac{M}{T}\grad n\right)
  +
  \frac{M_n|_T}{2}|\grad n|^2 .\sourceeqbadge{2.18}
  \label{eq:iib-muhat}
\end{equation}
The result is a variational derivative identity for the Section II.B
functionals.  It is not yet the reversible stress tensor, which is introduced
later in the source.

\subsection{Boundary Term and Its Role}

Substituting \eqref{eq:iib-muhat} into the integrated variation gives
\begin{equation}
  \delta S_b
  =
  \int_{\Omega}
  \left(
    \frac{1}{T}\delta\hat e
    -
    \frac{\hat\mu}{T}\delta n
  \right)\dd x
  -
  \int_{\partial\Omega}
  \frac{M}{T}(\nu\cdot\grad n)\delta n\,\dd a .\sourceeqbadge{2.20}
  \label{eq:iib-entropy-variation-boundary}
\end{equation}
This is the companion's reconstruction of \eqsrc{2.20}.  The sign comes from
the outward-normal convention in \eqref{eq:iib-ibp-identity}.  The boundary
integral is not optional bookkeeping: it determines which integrated variational
statement is being made.

Three common branches should be kept separate:

\begin{center}
\begin{tabular}{@{}p{0.22\linewidth}p{0.34\linewidth}p{0.34\linewidth}@{}}
\toprule
Branch & Local boundary term & What is still needed \\
\midrule
fixed density variation & vanishes because \(\delta n=0\) & no natural wall condition follows \\
natural boundary & coefficient of \(\delta n\) must vanish or be balanced & source wall data \\
global entropy/energy & boundary term is only one contribution & heat and wall-transfer closure \\
\bottomrule
\end{tabular}
\end{center}

Equivalently:
\begin{itemize}
\item If \(\delta n=0\) on \(\partial\Omega\), the displayed boundary term
  vanishes by the admissible-variation choice.
\item If \(\delta n\) is not fixed, a natural boundary condition or a wall
  entropy/energy functional must supply the missing surface contribution.
\item If a later global entropy or energy balance is claimed, this local
  Section II.B boundary term is not enough by itself; wall transfer and heat
  transfer data must also be specified.
\end{itemize}

\subsection{Coefficient Identities Used Later}

Because \(M\) depends on both \(n\) and \(T\), a recurring local identity is
\begin{equation}
  T\grad\left(\frac{M}{T}\right)
  =
  \grad M-\frac{M}{T}\grad T .
  \label{eq:iib-temperature-correction-identity}
\end{equation}
With \(M=CT+K\), this becomes
\begin{equation}
  \grad M-\frac{M}{T}\grad T
  =
  T\grad C+\grad K-\frac{K}{T}\grad T .
  \label{eq:iib-MCTK-temperature-correction}
\end{equation}
Under the simulation specialization later stated in the source, \(K=0\) and
\(C\) is constant, so the complete coefficient in
\eqref{eq:iib-MCTK-temperature-correction} vanishes.

For verification pedagogy, this companion also records the diagnostic
term-removal test expression \(\grad M\).  That test removes the
temperature-correction term in \eqref{eq:iib-temperature-correction-identity}
and asks what residual appears.  It is an audit-defined test expression, not a
source equation of Onuki or of Fukagawa.  Under \(K=0\) and constant \(C\), the
diagnostic test expression leaves \(C\grad T\), while the complete coefficient
vanishes.

\subsection{Type Checks}

\begin{itemize}
\item \(S_b\) is an integrated entropy; \(\hat S\) is an entropy density.
\item \(E_b\) is an integrated internal energy; \(\hat e\) is an
  internal-energy density including a gradient term.
\item \(M/T\) has the same variational role as \(C+K/T\), multiplying
  \(\grad n\cdot\grad\delta n\) in the entropy variation.
\item \(\hat\mu\) is a scalar per particle.  Its gradient is not introduced in
  this section as a force until the hydrodynamic equations are derived later.
\item The boundary term in \eqref{eq:iib-entropy-variation-boundary} is an
  entropy variation boundary term.  It is not a complete wall entropy balance.
\end{itemize}

\subsection{Verification Checklist}

The companion verification scripts for this section check only finite
identities:

\begin{enumerate}
\item the algebraic relation \(M=CT+K\);
\item the identity \(T\grad(M/T)=\grad M-(M/T)\grad T\);
\item the substitution
  \[
    \grad M-\frac{M}{T}\grad T
    =
    T\grad C+\grad K-\frac{K}{T}\grad T;
  \]
\item the \(K=0\), constant-\(C\) specialization of the complete coefficient;
\item the diagnostic term-removal test result \(C\grad T\) under that same
  specialization;
\item the one-dimensional version of the integration-by-parts identity
  \eqref{eq:iib-ibp-identity}.
\end{enumerate}

These checks support algebraic and variational bookkeeping.  They do not prove
the physical model, the second law, PDE well-posedness, boundary theory, or the
later numerical results.

The finite mirrors for this section are:
\begin{center}
\small
\begin{tabular}{@{}l@{}}
\texttt{scripts/python/section\_iib\_gradient\_entropy\_sympy.py}\\
\texttt{scripts/wolfram/section\_iib\_gradient\_entropy\_checks.wls}\\
\texttt{tests/test\_section\_iib\_gradient\_entropy.py}
\end{tabular}
\end{center}
Those scripts check the Eq.~(2.17) factor \(n\), the coefficient relation
\(M=CT+K\), the temperature-correction identity, and the
one-dimensional integration-by-parts residuals used above.


\subsection{Source-Procedure Closure for Section II.B}

The equation-inventory completion pass checks the source-procedure rows for
\eqsrc{2.12}--\eqsrc{2.20} against the derivations above.  The closure status is
bounded: it means that the local variational and coefficient steps are now
explicit in the companion, not that later hydrodynamic or wall-transfer branches
are closed.

\begin{center}
\begin{tabular}{@{}p{0.18\linewidth}p{0.28\linewidth}p{0.40\linewidth}@{}}
\toprule
Source item & Procedure type & Companion coverage \\
\midrule
\eqsrc{2.12}--\eqsrc{2.15} & functional definitions and type split & entropy, energy, and hatted densities typed before variation \\
\eqsrc{2.16}--\eqsrc{2.18} & fixed-\(\hat e\) variation & \(M=CT+K\), \(M_n|_T\), and \(\hat\mu\) derived before integration by parts \\
\eqsrc{2.20} & boundary variation & outward-normal boundary term displayed and branch table recorded \\
text after \eqsrc{2.19} & derivative convention & \(M_n|_T\) separated from the spatial derivative of \(M(n,T)\) \\
\bottomrule
\end{tabular}
\end{center}

The remaining open objects in this row group are not missing II.B derivations;
they are later-source objects: hydrodynamic stress, global wall transfer,
Appendix-B scaling, and unpublished simulation-code information.

\subsection*{Tagged verification complement}

\OnukiLinkIIB{The tagged Section II.B verification complement} records the
variation steps, executable mirrors, and the diagnostic negative control without
turning the local calculation into a global-transfer claim.

\subsection{Open Issues Carried Forward}

\begin{itemize}
\item The Appendix-B B3 printed-source branch versus dimensional-source branch
  issue is outside Section II.B and is deferred to the scaled-equation
  appendix.
\item Wall and global entropy-transfer closure are not settled by
  \eqsrc{2.20}; later wall functionals and heat-transfer terms are required.
\item The reversible stress tensor and its divergence are deferred to the
  hydrodynamic-equation section.
\item The admissible variation space at a physical wall must be stated before
  an integrated variational equivalence claim can be made.
\end{itemize}
